\documentclass[11pt]{article}
\usepackage[a4paper,margin=1in]{geometry}
\usepackage{amsmath,amssymb,amsthm,mathtools}
\usepackage{graphicx}
\usepackage[hidelinks]{hyperref}
\emergencystretch=3em

\newtheorem{theorem}{Theorem}
\newtheorem{lemma}[theorem]{Lemma}
\newtheorem{corollary}[theorem]{Corollary}
\newcommand{\D}{\mathcal D}
\newcommand{\rhoD}{\rho_{\mathcal D}}

\title{Prescribed-Coefficient Greedy Convergence for\\
Countable Orthogonal-Block Dictionaries}
\author{A substantial partial answer to the question in arXiv:2209.03091}
\date{August 11, 2026}

\begin{document}
\maketitle

\begin{abstract}
Oliaro--Tomatis--Valiullin--Valiullin ask for a characterization of the
symmetric Hilbert-space dictionaries for which every exact greedy expansion
converges for every positive coefficient sequence $c_n\to0$ with
$\sum_n c_n=\infty$.  We prove convergence for a new countable orthogonal-block
class.  Each block may be infinite-dimensional and carry an arbitrary
symmetric dictionary, provided all block dictionaries have a common positive
norming constant and attain their local maxima.  In particular, an
orthonormal basis may be augmented by arbitrary closed sets of atoms in
infinitely many pairwise disjoint finite coordinate blocks of uniformly
bounded size.  The source allowed extra atoms only in one fixed finite block.
\end{abstract}

\noindent\textbf{Status.}
\texttt{substantial partial result; likely valid}.  The full characterization
remains open.

\section{Source question and definitions}

Let $H$ be a real Hilbert space.  A symmetric dictionary $\D$ is a symmetric
set of unit vectors with dense linear span.  Put
\[
 \rhoD(x)=\sup_{g\in\D}\langle x,g\rangle
          =\sup_{g\in\D}|\langle x,g\rangle|.
\]
For positive prescribed coefficients $(c_n)$, the exact greedy algorithm starts
with $r_0=f$ and, while $r_{n-1}\ne0$, chooses $g_n\in\D$ satisfying
\[
 \langle r_{n-1},g_n\rangle=\rhoD(r_{n-1}),
 \qquad r_n=r_{n-1}-c_ng_n.
\]
Convergence means $r_n\to0$ for every initial $f$ and every possible exact
maximizer sequence.

The source proves this when $\D$ is a symmetrized orthonormal basis, and when
arbitrary additional atoms are supported on one fixed finite subset of that
basis.  Its concluding question asks for a characterization of all dictionaries
having convergence under
\begin{equation}\label{eq:coefficients}
 c_n>0,\qquad c_n\longrightarrow0,\qquad \sum_{n=1}^{\infty}c_n=\infty.
\end{equation}

\begin{center}
\includegraphics[width=.94\linewidth]{figures/open_question_crop.png}
\end{center}

\section{Proof intuition}

Write the Hilbert space as an orthogonal sum of blocks.  Greedy selection then
compares one scalar quantity from each block: its local dictionary support.
For the initial vector, order those local supports as
$s_1\ge s_2\ge\cdots$.  Before the first time the global greedy support falls
below $s_N$, no block after $N$ can have been touched.  At that stopping time,
the uniform local norming condition bounds each of the first $N$ block norms
by $s_N/\alpha$, while the remaining tail is still the initial tail.  Since
$(s_j)\in\ell_2$ and is decreasing, $Ns_N^2\to0$, forcing a subsequence of
remainders to tend to zero.

To upgrade subsequential decay to full convergence, we use an exact-greedy
Lyapunov quantity that seems particularly well suited to this question:
\[
 E_n=\|r_n\|^2-\rhoD(r_n)^2.
\]
Symmetry and exact maximization make $(E_n)$ nonincreasing.  Once
$\liminf\|r_n\|=0$, this defect tends to zero, so any remainder bounded away
from zero is almost parallel to a maximizing atom and its norm must decrease.

\section{A general exact-greedy lemma}

\begin{lemma}[Subsequential decay forces convergence]\label{lem:upgrade}
Let $\D$ be a symmetric dictionary for which the exact maximum is attained at
every remainder.  Suppose $c_n\to0$.  For any exact greedy trajectory,
\[
 \liminf_{n\to\infty}\|r_n\|=0
 \quad\Longrightarrow\quad
 \lim_{n\to\infty}\|r_n\|=0.
\]
\end{lemma}

\begin{proof}
Set $a_n=\rhoD(r_n)$ and $E_n=\|r_n\|^2-a_n^2\ge0$.  If $g_{n+1}$ is chosen
at step $n+1$, then $\langle r_n,g_{n+1}\rangle=a_n$.  Symmetry gives
\[
 a_{n+1}\ge
 |\langle r_n-c_{n+1}g_{n+1},g_{n+1}\rangle|
 =|a_n-c_{n+1}|.
\]
Consequently,
\begin{align*}
 E_{n+1}
 &\le \|r_n\|^2-2c_{n+1}a_n+c_{n+1}^2
       -(a_n-c_{n+1})^2\\
 &=E_n.
\end{align*}
Thus $E_n$ decreases.  The liminf hypothesis and $0\le E_n\le\|r_n\|^2$
force $E_n\to0$.

Fix $\varepsilon>0$.  On a sufficiently late tail,
$E_n\le\varepsilon^2/16$ and $c_{n+1}\le\varepsilon/4$.  Whenever
$\|r_n\|\ge\varepsilon$ on that tail,
\[
 a_n\ge\sqrt{\|r_n\|^2-E_n}\ge\frac{\sqrt{15}}4\varepsilon,
\]
and hence
\[
 \|r_{n+1}\|^2-\|r_n\|^2
 =-2c_{n+1}a_n+c_{n+1}^2<0.
\]
Choose a still later index $k$ with $\|r_k\|<\varepsilon$.  If a subsequent
remainder is below $\varepsilon$, the next has norm less than
$\varepsilon+c_{n+1}\le5\varepsilon/4$; while its norm is at least
$\varepsilon$, it decreases.  Therefore
$\sup_{n\ge k}\|r_n\|\le5\varepsilon/4$.  Since $\varepsilon$ is arbitrary,
$r_n\to0$.
\end{proof}

\section{Countable orthogonal-block theorem}

\begin{theorem}\label{thm:block}
Let
\[
 H=\bigoplus_{j=1}^{\infty}H_j
\]
be an orthogonal Hilbert sum.  For each $j$, let $\D_j$ be a symmetric unit
dictionary in $H_j$ such that every local support maximum
\[
 \rho_j(x)=\max_{g\in\D_j}\langle x,g\rangle
\]
is attained.  Suppose there is one $\alpha>0$ such that
\begin{equation}\label{eq:width}
 \rho_j(x)\ge\alpha\|x\|\qquad(x\in H_j,\ j\ge1).
\end{equation}
Embed the $\D_j$ into $H$ and put $\D=\bigcup_j\D_j$.  Then for every
coefficient sequence satisfying \eqref{eq:coefficients}, every exact greedy
expansion with respect to $\D$ converges to its initial vector.
\end{theorem}

\begin{proof}
For $x=(x_j)\in H$, the numbers $\rho_j(x_j)$ tend to zero, since they are at
most $\|x_j\|$.  Thus the global supremum is a maximum whenever $x\ne0$, and
$\D$ is an admissible symmetric dictionary.

Let $(r_n)$ be a greedy trajectory and write $a_n=\rhoD(r_n)$.  The energy
identity is
\begin{equation}\label{eq:energy}
 \|r_{n+1}\|^2
 =\|r_n\|^2-2c_{n+1}a_n+c_{n+1}^2.
\end{equation}
We first have $\liminf_n a_n=0$.  Otherwise $a_n\ge\delta>0$ eventually; once
$c_{n+1}\le\delta$, equation \eqref{eq:energy} gives
$\|r_{n+1}\|^2\le\|r_n\|^2-\delta c_{n+1}$.  The divergent sum in
\eqref{eq:coefficients} then forces a negative squared norm.

Write the initial vector as $f=(f_j)$ and relabel the blocks so that
\[
 s_j:=\rho_j(f_j),\qquad s_1\ge s_2\ge\cdots\ge0.
\]
We have $\sum_j s_j^2\le\sum_j\|f_j\|^2<\infty$.  If only finitely many
$s_j$ are nonzero, the trajectory lies in a finite orthogonal sum on which
\[
 \rhoD(x)=\max_j\rho_j(x_j)
 \ge\frac{\alpha}{\sqrt N}\|x\|,
\]
and the same energy argument as above gives convergence.  We may therefore
assume all $s_N>0$.

For each $N$, let $m_N$ be the first index with $a_{m_N}\le s_N$; it exists
because $\liminf a_n=0$.  Every block $j>N$ is untouched before this time.
Indeed, while it is untouched its local support is $s_j\le s_N$, whereas at
every earlier step the selected global support is strictly larger than $s_N$.
At time $m_N$, condition \eqref{eq:width} therefore gives
\[
 \|(r_{m_N})_j\|\le\frac{s_N}{\alpha}\quad(j\le N),
 \qquad (r_{m_N})_j=f_j\quad(j>N).
\]
Hence
\begin{equation}\label{eq:stop}
 \|r_{m_N}\|^2
 \le \frac{N s_N^2}{\alpha^2}+\sum_{j>N}\|f_j\|^2.
\end{equation}
Since $(s_j)$ is decreasing and square summable,
\[
 \frac N2s_N^2\le
 \sum_{j=\lceil N/2\rceil}^{N}s_j^2\longrightarrow0.
\]
The other term in \eqref{eq:stop} is the tail of the square-summable block
norms.  Thus $\|r_{m_N}\|\to0$.  Lemma \ref{lem:upgrade} completes the proof.
\end{proof}

\section{Infinitely many finite perturbation blocks}

\begin{corollary}\label{cor:coordinates}
Let $(e_i)_{i\in I}$ be an orthonormal basis and partition $I$ into finite
sets $I_j$ satisfying $|I_j|\le M<\infty$.  In each
$H_j=\operatorname{span}\{e_i:i\in I_j\}$, let $Y_j$ be any closed symmetric
set of unit vectors.  Then
\[
 \D=\{\pm e_i:i\in I\}\cup\bigcup_{j=1}^{\infty}Y_j
\]
has the convergence property of Theorem \ref{thm:block}.
\end{corollary}

\begin{proof}
For the local dictionary consisting of the signed coordinate vectors and
$Y_j$, compactness gives attainment, and
\[
 \rho_j(x)\ge\max_{i\in I_j}|\langle x,e_i\rangle|
 \ge\frac1{\sqrt M}\|x\|.
\]
Apply Theorem \ref{thm:block} with $\alpha=M^{-1/2}$.
\end{proof}

When infinitely many $Y_j$ are nonempty, their union is not supported on one
fixed finite subset of the basis.  Corollary \ref{cor:coordinates} is therefore
a strict extension of the source's Corollary 4.3.

\section{Scope, verification, and novelty status}

The theorem is limited to exact greediness $t=1$, orthogonal component
dictionaries, and a uniform positive local norming constant.  It does not
characterize nonorthogonal dictionaries such as general Riesz bases, and it
does not settle what happens when the best local constants tend to zero.
Those two upgrade routes were examined; Gram-matrix interactions destroy the
untouched-tail argument, while no rigorous counterexample was obtained in the
nonuniform block case.

The proof was audited through the two independent monotonic mechanisms
$\liminf a_n=0$ from \eqref{eq:energy} and monotonicity of
$\|r_n\|^2-a_n^2$ in Lemma \ref{lem:upgrade}.  A bounded search of the exact
title, arXiv id, characterization wording, and citing literature found the
2025 journal version of the source, which retains the open question, but no
later solution or countable-block theorem.  Human mathematical and literature
review is recommended.

\begin{thebibliography}{9}
\bibitem{OTVV}
A. Oliaro, L. Tomatis, A. R. Valiullin, and A. R. Valiullin,
\emph{Greedy expansions with prescribed coefficients in Hilbert spaces for
special classes of dictionaries}, Constructive Approximation 62 (2025),
243--259; arXiv:2209.03091.  See Corollary 4.3 and Section 5.
\end{thebibliography}

\end{document}
